\documentclass[conference]{IEEEtran}
\usepackage{cite}
\usepackage{amsmath,amssymb,amsfonts}
\usepackage{algorithmic}
\usepackage{graphicx}
\usepackage{textcomp}
\usepackage{xcolor}
\def\BibTeX{{\rm B\kern-.05em{\sc i\kern-.025em b}\kern-.08em
    T\kern-.1667em\lower.7ex\hbox{E}\kern-.125emX}}
\begin{document}

\title{Project: COVID-19 Assisted Diagnosis\\


\author{\IEEEauthorblockN{Carlo Heemeryck\textsuperscript{1}}
\and
\IEEEauthorblockN{Willem Van Nieuwenhuyse\textsuperscript{1}}
\and
\IEEEauthorblockN{Seppe Vanrietvelde\textsuperscript{2}}
\and
\IEEEauthorblockN{Luca Visser\textsuperscript{2}}
\and

\textsuperscript{1} Bachelor of Science in de ingenieurswetenschappen - werktuigkunde-elektrotechniek \hfill\\
\textsuperscript{2} Master of Science in Statistical Data Analysis - Computational Statistics\hfill\\}
}
\maketitle



\section{Introduction}
Adds context to your project (referencing the material in your bibliography is important~\cite{b1, b2}) and presents the content of the following sections (starting with Section~\ref{sec:task_1}).


\section{Task 1: Data Exploration, Pre-Processing and Augmentation}\label{sec:task_1}
You can freely organize the sections about the tasks. The given structure is purely a suggestion.


\subsection{Subsection 1 (e.g. Methods)}
 Details the strategies you have employed, including those that did not work.

\subsection{Subsection 2 (e.g. Results)}
Includes tables with quantitative results (Table~\ref{table:example}) and images (Fig.~\ref{fig:example}) from your project while carefully explaining their meaning and how you produced them.
\begin{table}[htbp]
\caption{Table Type Styles}
\begin{center}
\begin{tabular}{|c|c|c|c|}
\hline
\textbf{Table}&\multicolumn{3}{|c|}{\textbf{Table Column Head}} \\
\cline{2-4} 
\textbf{Head} & \textbf{\textit{Table column subhead}}& \textbf{\textit{Subhead}}& \textbf{\textit{Subhead}} \\
\hline
copy& More table copy$^{\mathrm{a}}$& &  \\
\hline
\multicolumn{4}{l}{$^{\mathrm{a}}$Sample of a Table footnote.}
\end{tabular}
\label{table:example}
\end{center}
\end{table}

\begin{figure}[htbp]
\centerline{\includegraphics{fig1.png}}
\caption{Example of a figure caption.}
\label{fig:example}
\end{figure}
\subsection{Subsection 3 (e.g. Discussion)}
Discusses the results and your general experience with the project, as well as provide answers to the questions about the task.
\section{Task 2: Building the baseline model}
\section{Task 3: Transfer Learning}
\section{Task 4: Explainability through Grad-CAM}
\section{Conclusions}
\section{Author contributions and collaboration}
\section{Use of Generative AI}


\subsection{Figures and Tables}
\paragraph{Positioning Figures and Tables}




\section*{References}
Example References:
\begin{thebibliography}{00}
\bibitem{b1} G. Eason, B. Noble, and I. N. Sneddon, ``On certain integrals of Lipschitz-Hankel type involving products of Bessel functions,'' Phil. Trans. Roy. Soc. London, vol. A247, pp. 529--551, April 1955.
\bibitem{b2} J. Clerk Maxwell, A Treatise on Electricity and Magnetism, 3rd ed., vol. 2. Oxford: Clarendon, 1892, pp.68--73.
\bibitem{b3} I. S. Jacobs and C. P. Bean, ``Fine particles, thin films and exchange anisotropy,'' in Magnetism, vol. III, G. T. Rado and H. Suhl, Eds. New York: Academic, 1963, pp. 271--350.
\bibitem{b4} K. Elissa, ``Title of paper if known,'' unpublished.
\bibitem{b5} R. Nicole, ``Title of paper with only first word capitalized,'' J. Name Stand. Abbrev., in press.
\bibitem{b6} Y. Yorozu, M. Hirano, K. Oka, and Y. Tagawa, ``Electron spectroscopy studies on magneto-optical media and plastic substrate interface,'' IEEE Transl. J. Magn. Japan, vol. 2, pp. 740--741, August 1987 [Digests 9th Annual Conf. Magnetics Japan, p. 301, 1982].
\bibitem{b7} M. Young, The Technical Writer's Handbook. Mill Valley, CA: University Science, 1989.
\end{thebibliography}
\end{document}
